\section{结论}

\subsection{主要结论}

本文针对圆柱形中药材热风烘干过程，建立了圆柱坐标下一维轴对称的耦合传热--传质模型，并以统一的数值框架求解了全部四个问题。主要结论如下。

\textbf{问题一（预热平衡阶段）}：该阶段热量尚未穿透整个断面而水分变化集中在表层。经计算，$1800$~s 时药材中心与表面温度分别为 \QOneCenterT~$^\circ$C 与 \QOneSurfaceT~$^\circ$C（相差 $3.21~^\circ$C，仍低于当时环境温度 $41.51~^\circ$C），中心与表面干基含水率分别为 \QOneCenterC~kg/kg 与 \QOneSurfaceC~kg/kg（表面已失去约 $41\%$ 的初始含水率，而中心几乎未变）。温度的作用范围在 $30$~min 内已覆盖整个半径，水分的作用范围仅为表面附近约 $0.15\Rzero$（穿透深度 $0.298$~cm），因此\textbf{不能以``含水率不变''为由忽略传质，也不能对全断面施加统一的失水修正}。

\textbf{问题二（整个烘干过程）}：采用附录 3 的变物性公式后，干燥过程的时间尺度完全由传质决定：温度场在约 $4$~h 内即与环境同温（$50.165~^\circ$C）并保持均匀，而含水率要经过 $2.38$ 天才能整体达标。全过程呈现三段式规律——前 $3$~h 表层快速失水（表面降至 $1.0081$~kg/kg）、$3\sim24$~h 干燥前沿向内部推进（轴心由 $1.77$ 降到 $0.2371$~kg/kg）、$24$~h 之后全场缓慢收敛。其机制是扩散系数随含水率下降而急剧减小：由初值处的 $1.36\times10^{-8}$~m$^2$/s 降到判据值处的 $8.05\times10^{-10}$~m$^2$/s，扩散特征时间由 $8.2$~h 增大到 $138$~h。整个烘干过程历时 \textbf{\QTwoEventH~h}。

\textbf{问题三（烘干时长）}：把``药材各处含水率低于 $0.15$~kg/kg''严格解释为全场最大值判据，并把它写成只依赖时间的事件函数后，求根问题化为一次积分中的向下过零检测。求得烘干所需时间为 \textbf{\QThreeEventH~h}（\QThreeEventS~s）；交付文件写到首个严格合格的 $60$~s 采样行 \QThreeLastRowS~s（\QThreeLastRowH~h），两者相差 \QThreeGap~s，含义不同、不可混用。结果还表明\textbf{烘干时长实际上由轴心附近最后 $0.5$~cm 的湿芯决定}：表面在 $12$~h 内即已达标，而轴心到 $57.17$~h 才达标，中间相隔约 $45$~h，占总时长的 $79\%$。

\textbf{问题四（考虑失水收缩）}：由干物质守恒导出干基表观密度按 ${\R}^{-2}$ 增大，并通过材料坐标变换 $\xi=\rpos/\R(\t)$ 证明了坐标变换引入的输运项与材料速度对流项\textbf{严格相消}，从而把移动边界问题化为固定域问题且方程中不含对流项。求得烘干时长为 \textbf{\QFourEventH~h}，结束时刻的药材半径为 \QFourEventR~cm（收缩 \QFourShrinkPct\%），且该时刻落在附件 2 半径平台的内部，故答案不依赖半径数据在终点附近的插值细节。经三个算例的对照分离后：参数集由附录 3 换为附录 4 使烘干时长按 $2.26$ 倍变慢，而\textbf{收缩本身使烘干时长缩短 \QFourShrinkEffectPct\%}（$129.09$~h $\to$ $50.82$~h）；两者叠加的净效应为缩短 \QFourNetEffectPct\%。收缩使干燥加快的机制是：在内部扩散控制下（$\hmass\R/\Ddiff\approx20\sim25\gg1$），时间尺度正比于 ${\R}^2/\Ddiff$，而干基表观密度随收缩上升抵消并超过了表面积减小的影响。

四问的最终答案汇总于表~\ref{tab:answers}。

\begin{table}[!ht]
  \centering
  \caption{四个问题的最终答案汇总}
  \label{tab:answers}
  \scriptsize
  \begin{tabular}{lll}
    \toprule
    问题 & 要求 & 本文答案 \\
    \midrule
    问题 1 & $1800$~s 内温度与含水率分布 & 中心/\allowbreak 表面温度 \QOneCenterT/\QOneSurfaceT~$^\circ$C；含水率 \QOneCenterC/\QOneSurfaceC~kg/kg \\
    问题 2 & 整个烘干过程的规律 & 温度 $4$~h 内均匀化；含水率三段式下降；全程 \QTwoEventH~h \\
    问题 3 & 烘干所需时间 & \textbf{\QThreeEventH~h}（\QThreeEventS~s） \\
    问题 4 & 考虑收缩的烘干时长 & \textbf{\QFourEventH~h}（\QFourEventS~s），末半径 \QFourEventR~cm \\
    \bottomrule
  \end{tabular}
\end{table}

\subsection{结论的适用范围}

上述结论成立的前提是第~\ref{sec:assumptions} 节列出的 15 条假设，其中三条对结论的适用范围影响最大，需要在使用这些数字时一并说明：

第一，\textbf{结论适用于热风对流为主导、且干燥由内部扩散控制的工况}。本文采用纯对流的热边界（不含蒸发耗热），并对传质使用题给的浓度差闭合。若实际过程以辐射或微波加热为主，或表面蒸发由外部气流控制而非内部扩散控制，则模型结构需要更换。

第二，\textbf{烘干时长的数值带有小时量级的不确定性}。$4$~h 之后的环境条件按末值保持外推覆盖了全程 $93\%$ 的时长，按附件 1 末小时的实测波动幅度做灵敏度分析，问题三的时长在 $57.04\sim57.91$~h 之间变化（$\pm1.3\%$）；而传质基准与潜热缺失这两项结构性闭合问题的影响无法用题目数据量化，方向上会使干燥偏快。因此``约 \QThreeEventH~h''应理解为当前模型与假设下的预测值。

第三，\textbf{问题四关于收缩效应的方向性结论依赖于传质闭合形式}。本文证明在本模型采用的干基通量闭合下，收缩使干燥加快，机制与 ${\R}^2/\Ddiff$ 一致；若改用按表面积定义的蒸气通量闭合，面积减小的减速效应会重新出现。因此该结论应表述为``在本模型闭合下收缩使干燥加快''，而不宜作为一般性物理结论。

在明确以上边界后，本文的模型与结果可用于：估算给定工艺条件下的烘干时间与所需批次间隔；判断工艺瓶颈在表面换热还是内部扩散，从而指导工艺参数的调整方向（本文的分析表明瓶颈在内部扩散，故提高温度比提高风速更有效）；以及在扩展为优化问题后，用于设计升温曲线以缩短烘干时间或降低能耗。
